\documentclass[12pt, a4paper]{article}

\usepackage{amsmath, amssymb, amsthm}
\usepackage{mathtools}
\usepackage{geometry}
\usepackage{hyperref}
\usepackage{booktabs}
\usepackage{array}
\usepackage{enumitem}
\usepackage{titlesec}
\usepackage{setspace}
\usepackage{fancyhdr}
\usepackage{natbib}

\geometry{top=1.2in, bottom=1.2in, left=1.3in, right=1.3in}
\setlength{\headheight}{14.5pt}

\hypersetup{colorlinks=true, linkcolor=black, citecolor=black, urlcolor=black}

\pagestyle{fancy}
\fancyhf{}
\renewcommand{\headrulewidth}{0.4pt}
\fancyhead[L]{\small\textit{Distinguishability Geometry}}
\fancyhead[R]{\small\thepage}

\titleformat{\section}{\large\bfseries}{\thesection.}{0.6em}{}
\titleformat{\subsection}{\normalsize\bfseries}{\thesubsection.}{0.6em}{}
\titleformat{\subsubsection}{\normalsize\itshape}{\thesubsubsection.}{0.6em}{}

\theoremstyle{plain}
\newtheorem{theorem}{Theorem}[section]
\newtheorem{proposition}[theorem]{Proposition}
\newtheorem{corollary}[theorem]{Corollary}
\newtheorem{lemma}[theorem]{Lemma}

\theoremstyle{definition}
\newtheorem{definition}[theorem]{Definition}
\newtheorem{example}[theorem]{Example}
\newtheorem{conjecture}[theorem]{Conjecture}

\theoremstyle{remark}
\newtheorem{remark}[theorem]{Remark}
\newtheorem{principle}[theorem]{Principle}

\newcommand{\deficit}{\delta_T}
\newcommand{\Dstar}{D^*}
\newcommand{\DT}{D_T}
\newcommand{\Lreach}{L_T}
\newcommand{\Lopt}{L^*}
\newcommand{\simX}{{\sim}_X}
\newcommand{\simY}{{\sim}_Y}
\newcommand{\CLIO}{\textsc{clio}}
\newcommand{\RSVP}{\textsc{rsvp}}
\newcommand{\Ddelta}{\Delta\delta}
\newcommand{\eps}{\epsilon}

\setstretch{1.15}
\setlength{\parskip}{0.4em}
\setlength{\parindent}{1.4em}

\begin{document}

% ── Title page ───────────────────────────────────────────────────────────────
\begin{titlepage}
\vspace*{2cm}
\begin{center}
{\LARGE\bfseries Distinguishability Geometry:\\[0.4em]
Projection, Embedding, Revision, and Transport}

\vspace{1.8cm}
{\large Flyxion}\\[0.4em]
{\normalsize Independent Researcher}

\vspace{1.2cm}
{\normalsize June 2026}

\vspace{2.2cm}
\begin{minipage}{0.82\textwidth}
\begin{abstract}
\noindent
We propose \emph{distinguishability geometry} as a unified framework for
understanding how representational systems gain, lose, and transfer the
capacity to express distinctions. The foundational primitive is the
\emph{ontological triple} $(X,\sim,T)$: a set $X$, an observational
indistinguishability relation $\sim$, and a generating ontology $T$. Every
such triple carries two intrinsic invariants: the \emph{coding deficit}
$\delta_T(x)=L_T(x)-L^*(x)$, measuring excess description length, and the
\emph{distinguishability deficit} $\Delta_T(x)=D^*(x)-D_T(x)$, measuring
inaccessible distinction capacity. The \emph{Deficit Proxy Theorem} establishes
$\Delta_T(x)\le C\,\delta_T(x)$, making the coding deficit a computable
surrogate for the geometric one; full equivalence is conjectured but not
proved.

We identify four primitive operations arranged in a two-level hierarchy.
At the \emph{distinction level}, projection and embedding coarsen and
refine $\sim$ within a fixed space. At the \emph{ontology level}, revision
restructures the generating ontology and transport maps distinguishability
structure across spaces. The \emph{Representational Reduction Proposition}
establishes these as exhaustive qualitative fiber effects; the
\emph{Classification Conjecture} asserts every representational
transformation factors into their finite composition.

For each operation we derive a characteristic theorem: a Projection Deficit
Bound in entropic form, an Embedding Refinement theorem, a Revision
Monotonicity theorem with asymptotic corollary, and a Transport Stability
theorem giving a Lipschitz condition for analogy. A \emph{Distortion--Deficit
Relationship} connects the two error measures. A concrete finite example
illustrates all four operations. We close with a categorical formulation in
which the hierarchy corresponds to morphism types and the deficit is a functor
from the category $\mathbf{Dist}$ of ontological triples.
\end{abstract}
\end{minipage}
\end{center}
\end{titlepage}

\tableofcontents
\newpage

% ─────────────────────────────────────────────────────────────────────────────
\section{Introduction}
\label{sec:intro}
% ─────────────────────────────────────────────────────────────────────────────

Three independent research programmes have converged on a common mathematical
object. The first, from process-primary computation, establishes that observable
state is a projection of a richer history space and that distinct histories
collapse to the same state. The second, from kernel statistics, demonstrates
that distinct distributions become mutually singular after covariance embedding,
amplifying separation without creating information. The third, from archive-based
compression, introduces \emph{reachable complexity}: the shortest description
expressible from the archive's current template hierarchy, which may be strictly
longer than the Kolmogorov-optimal description.

Each programme arrives at the same hidden structure from a different direction:
the \emph{geometry of distinguishability}. The first studies distinctions
\emph{destroyed} by representation. The second studies distinctions
\emph{amplified}. The third studies distinctions that are \emph{inaccessible}
within the current frame.

The present paper makes this object explicit. Our main contributions are:

\begin{enumerate}[leftmargin=1.8em, label=(\arabic*)]
  \item The \emph{ontological triple} $(X,\sim,T)$ as the correct object of
        study, fixing the under-specification of earlier formulations that left
        the generating ontology implicit.
  \item Two intrinsic invariants---coding deficit and distinguishability
        deficit---and a \emph{Deficit Proxy Theorem} relating them.
  \item A two-level hierarchy of four operations with characteristic theorems
        for each, including a corrected Embedding Refinement proof, an entropic
        Projection Deficit Bound, and a Lipschitz Transport Stability theorem.
  \item A \emph{Distortion--Deficit Relationship} connecting the two error
        measures that previously lived in separate worlds.
  \item A \emph{Classification Conjecture} precisely stated as a factorization
        claim, together with a worked finite example illustrating all four
        operations.
  \item A categorical formulation of $\mathbf{Dist}$ with well-defined objects
        and morphism types.
\end{enumerate}

The paper is organized as: spaces and invariants (\S\ref{sec:spaces}),
classification conjecture (\S\ref{sec:conjecture}), four operations with
theorems (\S\ref{sec:operations}), asymmetries and meta-theorem
(\S\ref{sec:asymmetry}), deficit as universal invariant
(\S\ref{sec:deficit-universal}), and categorical formulation
(\S\ref{sec:category}).

% ─────────────────────────────────────────────────────────────────────────────
\section{Distinguishability Spaces and Their Invariants}
\label{sec:spaces}
% ─────────────────────────────────────────────────────────────────────────────

\subsection{The Ontological Triple}

Earlier formulations defined the foundational object as a pair $(X,\sim)$.
This is insufficient: the indistinguishability relation $\sim$ on $X$ does
not by itself determine which descriptions are reachable, and hence does
not determine the deficit. The correct object is a triple.

\begin{definition}[Ontological triple]
\label{def:triple}
An \emph{ontological triple} is $(X,\sim,T)$ where $X$ is a set, $\sim$ is
an equivalence relation on $X$ (observational indistinguishability), and
$T$ is an ontology: a description language or template hierarchy determining
which descriptions of elements of $X$ are reachable. The equivalence classes
$[x]_\sim$ are called \emph{fibers}.
\end{definition}

The same pair $(X,\sim)$ can appear in multiple triples with different ontologies
$T$, $T'$, inducing the same fiber structure but different deficit values.
This is not a defect; it reflects the genuine dependence of representational
cost on the description language.

\begin{proposition}[Fiber Monotonicity]
\label{prop:monotone}
Let $(X,\sim_1,T)$ and $(X,\sim_2,T)$ share the same set and ontology with
$\sim_1\subseteq\sim_2$ ($\sim_2$ coarser). Then:
\[
  [x]_{\sim_1}\subseteq[x]_{\sim_2},\qquad
  |[x]_{\sim_1}|\le|[x]_{\sim_2}|.
\]
\end{proposition}

\begin{proof}
If $y\in[x]_{\sim_1}$ then $x\sim_1 y$; since $\sim_1\subseteq\sim_2$,
$x\sim_2 y$, so $y\in[x]_{\sim_2}$. Cardinality follows.
\end{proof}

\begin{corollary}
Coarsening a distinguishability relation never increases the number of
distinguishable classes. Projection always reduces expressible distinctions.
\end{corollary}

\subsection{The Distinguishability Principle}

\begin{principle}[Distinguishability Principle]
\label{prin:dist}
A representational system is characterized not by what objects it contains
but by what distinctions it can express.
\end{principle}

Every ontological triple possesses a deficit because no finitary ontology
can express all distinctions present in the world.

\subsection{Two Deficit Notions}
\label{subsec:deficits}

\begin{definition}[Coding deficit]
\label{def:coding-deficit}
Let $\Lreach(x)$ be the shortest description of $x$ reachable within $T$,
and $\Lopt(x)=K(x)$ the Kolmogorov complexity of $x$. The \emph{coding
deficit} is:
\[
  \delta_T(x) \;=\; \Lreach(x) - \Lopt(x).
\]
\end{definition}

\begin{definition}[Distinguishability deficit]
\label{def:dist-deficit}
Let $D_T(x)$ be the distinguishability capacity of ontology $T$ at $x$:
the number of objects separable from $x$ by descriptions reachable in $T$.
Let $D^*(x)$ be the maximum capacity attainable by any ontology. The
\emph{distinguishability deficit} is:
\[
  \Delta_T(x) \;=\; D^*(x) - D_T(x).
\]
\end{definition}

\begin{theorem}[Deficit Non-Negativity]
\label{thm:nonneg}
For every ontological triple $(X,\sim,T)$ and every $x\in X$:
\[
  \delta_T(x)\ge 0,\qquad \Delta_T(x)\ge 0.
\]
Both vanish iff $T$ is locally optimal for $x$.
\end{theorem}

\begin{proof}
$\Lopt(x)$ is the minimum over all ontologies, so $\Lopt(x)\le\Lreach(x)$,
giving $\delta_T(x)\ge 0$. Similarly $D^*(x)$ is the maximum, so
$D_T(x)\le D^*(x)$, giving $\Delta_T(x)\ge 0$. Both equalities hold iff
$T$ achieves the respective optimum.
\end{proof}

\begin{theorem}[Deficit Proxy]
\label{thm:proxy}
Under the assumption that each expressible distinction requires at least one
description bit (\emph{faithful encoding}), the coding deficit bounds the
distinguishability deficit:
\[
  \Delta_T(x) \;\le\; C\,\delta_T(x)
\]
for a system-dependent constant $C>0$. In particular:
\[
  \delta_T(x)=0 \;\Longrightarrow\; \Delta_T(x)=0.
\]
The converse direction, $\Delta_T(x)=0\Rightarrow\delta_T(x)=0$, requires
the additional assumption that optimal coding uses all available distinctions,
and is left as a conjecture.
\end{theorem}

\begin{proof}
Under faithful encoding, each inaccessible distinction in $\Delta_T(x)$
requires at least one additional bit in $T$'s description of $x$ to flag
the corresponding fiber collapse. Hence $\Delta_T(x)\le\delta_T(x)$,
and the general bound $\Delta_T(x)\le C\,\delta_T(x)$ follows with $C\ge 1$
scaling with the maximum description-length cost per inaccessible distinction.

For the forward implication: if $\delta_T(x)=0$, then $T$ achieves optimal
description length for $x$. Under faithful encoding, this means every
bit-costly distinction is expressible in $T$, so no distinction is
inaccessible: $\Delta_T(x)=0$.
\end{proof}

\begin{conjecture}[Full Deficit Correspondence]
\label{conj:full-corr}
Under faithful encoding, $\Delta_T(x)=0\iff\delta_T(x)=0$.
\end{conjecture}

\begin{remark}
The Deficit Proxy Theorem is sufficient for all subsequent arguments in
this paper. The coding deficit $\delta_T(x)$ is a computable surrogate
for the geometric quantity $\Delta_T(x)$, and the two share the same zero
set under the proxy direction. Full equivalence remains an open problem.
\end{remark}

\subsection{A Concrete Illustration}
\label{subsec:example}

Before stating the classification conjecture, we ground the framework in a
finite example that will be extended throughout the paper.

\begin{example}[Four operations on a four-element set]
\label{ex:four}
Let $X=\{a,b,c,d\}$ with ontology $T_0$ admitting a full distinguishing
description for each element. Consider the following relations and ontologies:

\medskip
\noindent\textbf{Initial space.}
$\sim_0 = \bigl\{\{a\},\{b\},\{c\},\{d\}\bigr\}$ (discrete: all elements
distinguishable). Here $\delta_{T_0}(x)=0$ for all $x$.

\medskip
\noindent\textbf{Projection.}
The map $\sigma:X\to\{p,q,r\}$ with $\sigma(a)=\sigma(b)=p$, $\sigma(c)=q$,
$\sigma(d)=r$ induces $\sim_1=\bigl\{\{a,b\},\{c\},\{d\}\bigr\}$. The
fiber $F_p=\{a,b\}$ has size 2, so the projection deficit bound gives
$\Ddelta(\sigma)\ge\log 2=1$ bit.

\medskip
\noindent\textbf{Embedding.}
From $(X,\sim_0)$, map $\phi:X\to X\times\{0,1\}$ by $\phi(x)=(x,f(x))$
where $f(a)=f(c)=0$, $f(b)=f(d)=1$. This refines $\sim_0$ by separating
elements according to a new attribute: $\phi(a)\not\sim'\phi(b)$ even when
$a\sim_0 b$ under a coarser version of $\sim_0$. A faithful embedding
here preserves all existing distinctions and adds new ones.

\medskip
\noindent\textbf{Revision.}
Starting from $(X,\sim_1,T_0)$ with $\sim_1=\bigl\{\{a,b\},\{c\},\{d\}\bigr\}$,
suppose a new ontology $T_1\supset T_0$ adds a template that distinguishes
$a$ from $b$ directly. Then $\sim_{T_1}=\sim_0$ is recovered and
$\delta_{T_1}(a)=\delta_{T_1}(b)=0$: a productive revision has decreased
the deficit.

\medskip
\noindent\textbf{Transport.}
Let $Y=\{w,x,y,z\}$ with $\sim_Y=\bigl\{\{w,x\},\{y\},\{z\}\bigr\}$.
The map $\tau:X\to Y$ given by $\tau(a)=w$, $\tau(b)=x$, $\tau(c)=y$,
$\tau(d)=z$ is a transport with zero distortion: it preserves the fiber
structure of $\sim_1$ exactly ($a\sim_1 b\iff\tau(a)\sim_Y\tau(b)$).

\medskip
\noindent\textbf{Hybrid operation.}
Note that the transformation $\sim_0\to\sim_3$ where
$\sim_3=\bigl\{\{a,b\},\{c,d\}\bigr\}$ is \emph{not} a pure projection
(it coarsens in two places simultaneously). By the Classification Conjecture
(\S\ref{sec:conjecture}), it should factor as a composition: e.g.\ first
project $\{c,d\}$ to a single class, then project $\{a,b\}$. This factorization
is not unique, but the conjecture asserts it always exists.
\end{example}

% ─────────────────────────────────────────────────────────────────────────────
\section{The Classification Conjecture}
\label{sec:conjecture}
% ─────────────────────────────────────────────────────────────────────────────

Before developing the four operations in detail, we state the conjecture
that organizes the entire framework. The preceding example and the
Representational Reduction Proposition (stated as Proposition~\ref{prop:rep-reduction}
in \S\ref{sec:operations}) motivate it; the following sections constitute
the evidence for it.

\begin{conjecture}[Classification Conjecture]
\label{conj:classification}
Every representational transformation factors into a finite composition of
projection, embedding, revision, and transport maps.
\end{conjecture}

Three clarifications are essential.

\medskip
\noindent\textbf{Factorization, not identity.}
The transformation itself need not be one of the four primitive operations.
A hybrid transformation that simultaneously coarsens some fibers and refines
others factors as a composition---for instance, a projection followed by an
embedding. This is precisely the sense in which the conjecture resembles
factorization theorems in algebra and topology.

\medskip
\noindent\textbf{Scope.}
The conjecture concerns representational transformations: maps between
ontological triples $(X,\sim,T)$ that may alter any or all of $X$, $\sim$,
and $T$. Pure set-theoretic maps that ignore the relational and ontological
structure are outside its scope.

\medskip
\noindent\textbf{Evidence.}
The characteristic theorems of \S\ref{sec:operations} can be read as
partial evidence: they show that the four operations, together, are capable
of expressing the full range of qualitative fiber behaviors. Whether every
transformation \emph{factors} into these behaviors---whether no genuinely
new operation can arise from their combination---is what the conjecture
claims and what remains to be proved.

% ─────────────────────────────────────────────────────────────────────────────
\section{The Four Operations}
\label{sec:operations}
% ─────────────────────────────────────────────────────────────────────────────

\subsection{The Two-Level Structure}

The four operations divide into two levels according to what they modify:

\[
  \text{Objects} \;\longleftrightarrow\; \text{Distinctions}
  \;\longleftrightarrow\; \text{Ontologies}.
\]

Projection and embedding operate at the \emph{distinction level}: they
modify $\sim$ within a fixed space, moving horizontally between finer and
coarser relations on the same $X$. Revision and transport operate at the
\emph{ontology level}: revision changes $T$ (moving vertically in the diagram)
and transport connects two distinct triples (moving between nodes).

\begin{proposition}[Representational Reduction]
\label{prop:rep-reduction}
Every representational transformation induces a transformation on
distinguishability structure. The qualitatively distinct fiber effects are:
\begin{itemize}[leftmargin=1.8em]
  \item \textbf{Projection}: merges fibers (coarsens $\sim$).
  \item \textbf{Embedding}: separates fibers (refines $\sim$).
  \item \textbf{Revision}: restructures the ontology generating $\sim$.
  \item \textbf{Transport}: maps fiber structure between distinct triples.
\end{itemize}
These four effects are mutually exclusive as pure operations and jointly
exhaust the qualitative possibilities.
\end{proposition}

\begin{table}[h]
\centering
\renewcommand{\arraystretch}{1.35}
\begin{tabular}{lllll}
\toprule
\textbf{Level} & \textbf{Operation} & \textbf{Effect on $\sim$}
  & $\Ddelta$ & \textbf{Failure mode} \\
\midrule
Distinction & Projection  & Coarsens $\sim$           & $\ge 0$ & Information loss \\
Distinction & Embedding   & Refines $\sim$            & $\le 0$ & Spurious distinctions \\
Ontology    & Revision    & Restructures $T$          & $<0$ (productive) & Lock-in \\
Ontology    & Transport   & Maps $\sim$ across spaces & bounded & Fiber distortion \\
\bottomrule
\end{tabular}
\caption{The four operations, their level, deficit change $\Ddelta(f)$,
and canonical failure modes. Note: $\eps(f)\ge 0$ always; $\Ddelta(f)$
is signed.}
\label{tab:ops}
\end{table}

\subsection{Formal Separation: Distortion vs.\ Deficit Change}
\label{subsec:formalism}

\begin{definition}[Distortion and deficit change]
For a map $f$ between ontological triples:
\begin{itemize}[leftmargin=1.8em]
  \item $\eps(f)\ge 0$: the \emph{distortion} of $f$, measuring how much $f$
        fails to be a fiber-isomorphism. Always non-negative.
  \item $\Ddelta(f)=\delta_{T'}(f(x))-\delta_T(x)$: the \emph{deficit change}
        induced by $f$. Signed: positive means deficit increases, negative
        means it decreases.
\end{itemize}
\end{definition}

These quantities are related. The following theorem provides the connection
that was previously missing.

\begin{theorem}[Distortion--Deficit Relationship]
\label{thm:dist-deficit}
Let $f:(X,\sim_X,T_X)\to(Y,\sim_Y,T_Y)$ be a representational transformation
with distortion $\eps(f)$. Then under faithful encoding:
\[
  |\Ddelta(f)| \;\le\; K\,\eps(f)
\]
for a constant $K>0$ depending on the description complexity of individual
fibers in $T_X$ and $T_Y$.
\end{theorem}

\begin{proof}
Each unit of distortion $\eps(f)$ corresponds to a fiber violation: either
a pair that was equivalent in $X$ becoming inequivalent in $Y$ (distinction
inflation) or a pair inequivalent in $X$ becoming equivalent in $Y$
(distinction collapse). Under faithful encoding, each such violation costs
at most $K$ bits of description length: collapsing a fiber saves at most
$\log|F|$ bits (gained by the projection) and inflating creates at most
$\log|F'|$ bits (cost of the new separation). In both cases, the change in
deficit is bounded by the description complexity of the affected fibers,
which is controlled by $K\cdot\eps(f)$ where $K$ is the maximum
per-violation description cost.
\end{proof}

\begin{remark}
Theorem~\ref{thm:dist-deficit} makes distortion and deficit change into
\emph{exchange rates} between geometry and information. A transformation
with small distortion cannot change the deficit by much; a large deficit
change must be accompanied by substantial distortion. This connects the
geometric and coding-theoretic aspects of the framework into a single
inequality.
\end{remark}

\subsection{Projection}

\begin{definition}[Projection]
A \emph{projection} on $(X,\sim,T)$ is a surjective map $\sigma:X\to S$
inducing a coarsening $\sim'$ of $\sim$: $\sigma(x)=\sigma(y)\implies x\sim'y$.
Fibers under $\sim'$ are unions of fibers under $\sim$.
\end{definition}

\begin{theorem}[Projection Deficit Bound]
\label{thm:proj-deficit}
Let $\sigma:X\to S$ be a projection inducing fiber partition
$\{F_s=\sigma^{-1}(s)\}_{s\in S}$. Let $X$ carry a probability measure $\mu$.
Then the deficit increase satisfies:
\[
  \Ddelta(\sigma) \;\ge\; H(X\mid S)
  \;=\; -\sum_{s\in S}\mu(F_s)\sum_{x\in F_s}\frac{\mu(\{x\})}{\mu(F_s)}
  \log\frac{\mu(\{x\})}{\mu(F_s)},
\]
the conditional entropy of $X$ given $S$.
\end{theorem}

\begin{proof}
Describing $x\in X$ via $\sigma(x)\in S$ plus a within-fiber index requires
at least $H(X\mid S)$ additional bits beyond the description of $\sigma(x)$,
by Shannon's source coding theorem. Projection discards this index, making
$H(X\mid S)$ bits of distinguishing information inaccessible. Hence
$\Ddelta(\sigma)\ge H(X\mid S)$.
\end{proof}

\begin{corollary}
When fibers are uniform of size $|F_s|$, the bound reduces to
$\Ddelta(\sigma)\ge\log|F_s|$, recovering the earlier finite-fiber result
as a special case.
\end{corollary}

The conservation law governing projection is:
\begin{equation}
  |H_t| + |\Omega_t| = |\Omega_0|,
  \label{eq:conservation}
\end{equation}
where $H_t$ is the realized history, $\Omega_t$ the remaining possibility
space, and $|\cdot|$ measures information content. Projection transfers
possibility from the accessible history into the inaccessible remainder.

\subsection{Embedding}

\begin{definition}[Embedding]
An \emph{embedding} of $(X,\sim,T)$ into $(Y,\sim',T')$ is an injective map
$\phi:X\to Y$ such that the fiber structure is preserved:
$\phi(x)\sim'\phi(y)\implies x\sim y$. An embedding is \emph{faithful} if
the converse also holds: $x\sim y\iff\phi(x)\sim'\phi(y)$.
\end{definition}

\begin{theorem}[Embedding Refinement]
\label{thm:embed-refine}
Let $\phi:(X,\sim,T)\to(Y,\sim',T')$ be a faithful embedding. Then:
\[
  \delta_{T'}(\phi(x)) \;\le\; \delta_T(x).
\]
A faithful embedding never increases the coding deficit.
\end{theorem}

\begin{proof}
We need to show $L_{T'}(\phi(x))-L^*(\phi(x))\le L_T(x)-L^*(x)$.

Since $\phi$ is faithful, every description of $x$ reachable in $T$ corresponds
to a description of $\phi(x)$ reachable in $T'$: the image $\phi$ preserves
the full fiber structure, so $L_{T'}(\phi(x))\le L_T(x)$.

For the optimal lengths: $\phi$ is an injective map, so any optimal description
of $\phi(x)$ can be decoded to an optimal description of $x$ by applying
$\phi^{-1}$ to the fiber structure. This means $L^*(\phi(x))\ge L^*(x)$:
the optimal description of $\phi(x)$ cannot be shorter than the optimal
description of $x$, because the description must encode at least as much
distinguishing information (the embedding may add structure in $Y$ that is
extrinsic to $x$, making optimal descriptions of $\phi(x)$ longer, not shorter).

Combining: $\delta_{T'}(\phi(x))=L_{T'}(\phi(x))-L^*(\phi(x))
\le L_T(x)-L^*(x)=\delta_T(x)$.
\end{proof}

\begin{remark}
The key inequality is $L^*(\phi(x))\ge L^*(x)$, not $\le$. An injective
embedding into a richer space means a description of the image $\phi(x)$
must specify $\phi(x)$'s position within $Y$, which contains $X$ as a proper
subset. This can only make optimal descriptions longer (or equal), not shorter.
Non-faithful embeddings that collapse fibers may admit shorter optimal
descriptions for the image, but then they are not embeddings in the present
sense.
\end{remark}

\subsection{Revision}

\begin{definition}[Revision]
An \emph{ontological revision} is a transition $(X,\sim_T,T)\to(X,\sim_{T'},T')$
where $T'\supset T$ (every description reachable in $T$ is reachable in $T'$).
A revision is \emph{productive} if $\delta_{T'}(x)<\delta_T(x)$.
\end{definition}

\begin{theorem}[Revision Monotonicity]
\label{thm:revision-mono}
Let $T_0\subset T_1\subset T_2\subset\cdots$ be an expanding sequence of
ontologies. Then for all $n\ge 0$:
\[
  \delta_{T_{n+1}}(x) \;\le\; \delta_{T_n}(x).
\]
\end{theorem}

\begin{proof}
Since $T_n\subset T_{n+1}$, every description reachable in $T_n$ is reachable
in $T_{n+1}$, so $L_{T_{n+1}}(x)\le L_{T_n}(x)$. The optimal length
$L^*(x)=K(x)$ is independent of the ontology. Hence
$\delta_{T_{n+1}}(x)=L_{T_{n+1}}(x)-K(x)\le L_{T_n}(x)-K(x)=\delta_{T_n}(x)$.
\end{proof}

\begin{corollary}[Asymptotic deficit]
\label{cor:asymptotic}
The sequence $(\delta_{T_n}(x))_{n\ge 0}$ is monotone non-increasing and
bounded below by zero. Therefore:
\[
  \lim_{n\to\infty}\delta_{T_n}(x) \;=\; \delta_\infty(x) \;\ge\; 0.
\]
The \emph{residual deficit} $\delta_\infty(x)$ exists for every revision
sequence. $\delta_\infty(x)=0$ iff there exists a productive revision
sequence converging to an optimal ontology for $x$.
\end{corollary}

The canonical failure of revision is \emph{ontological lock-in}: $\delta_T(x)$
is large and $\delta_\infty(x)>0$, meaning no revision sequence reachable
from within $T$ can achieve optimality for $x$.

\subsection{Transport}

\begin{definition}[Transport]
\label{def:transport}
A \emph{transport map} is a function
$\tau:(X,\sim_X,T_X)\to(Y,\sim_Y,T_Y)$ with distortion at most $\eps\ge 0$:
\[
  x\sim_X x' \;\implies\; \tau(x)\sim_Y\tau(x'),
\]
and the measure of pairs $(\tau(x),\tau(x'))$ with $\tau(x)\sim_Y\tau(x')$
but $x\not\sim_X x'$ is controlled by $\eps$.
\end{definition}

\begin{theorem}[Transport Stability]
\label{thm:transport-stability}
Let $\tau:(X,\sim_X,T_X)\to(Y,\sim_Y,T_Y)$ be a transport map with
distortion at most $\eps$. Let $I_X:X\to\mathbb{R}$ be an invariant constant
on fibers of $\sim_X$, and let $I_Y:Y\to\mathbb{R}$ be a corresponding
invariant constant on fibers of $\sim_Y$ with $I_Y(\tau(x))$ well-defined
for all $x\in X$. Suppose $I_X$ and $I_Y$ agree on perfectly transported
fibers: $I_X(x)=I_Y(\tau(x))$ whenever $\tau$ maps the fiber of $x$ in $X$
isomorphically to the fiber of $\tau(x)$ in $Y$. Then:
\[
  |I_X(x) - I_Y(\tau(x))| \;\le\; C\eps
\]
for a constant $C$ depending on the oscillation of $I_X$ and $I_Y$ across
fiber boundaries.
\end{theorem}

\begin{proof}
Let $E_\eps\subseteq X$ be the set of $x$ for which $\tau$ violates
fiber preservation (either collapsing a distinction or creating a spurious
one). By the distortion bound, $\mu(E_\eps)\le\eps$ for some natural
measure $\mu$ on $X$. For $x\notin E_\eps$, the fiber of $x$ maps
isomorphically to the fiber of $\tau(x)$, so $I_X(x)=I_Y(\tau(x))$ by
assumption.

For $x\in E_\eps$, the invariants can differ by at most the oscillation of
$I_X$ (or $I_Y$) over a single fiber: $|I_X(x)-I_Y(\tau(x))|\le\mathrm{osc}(I)$.
Taking $C=\mathrm{osc}(I)/\min(\mu(F))$ where the minimum is over non-trivial
fibers, the bound $|I_X(x)-I_Y(\tau(x))|\le C\eps$ holds uniformly.
\end{proof}

\begin{corollary}[Justification of analogy]
An analogical inference transported by $\tau$ is valid up to error $C\eps(\tau)$.
An analogy is \emph{justified} for inference purposes iff $\eps(\tau)$ is
controlled and the relevant invariants have bounded oscillation. Transport
Stability is the precise mathematical condition under which analogical reasoning
is valid.
\end{corollary}

% ─────────────────────────────────────────────────────────────────────────────
\section{Asymmetries, Failure Modes, and the Meta-Theorem}
\label{sec:asymmetry}
% ─────────────────────────────────────────────────────────────────────────────

\subsection{Projection and Embedding Are Not Duals}

Both are distinction-level operations, but they are not inverses. Projection
necessarily has $\Ddelta(\sigma)\ge H(X\mid S)>0$ for non-trivial fibers:
it always loses information. Embedding has $\Ddelta(\phi)\le 0$ for faithful
embeddings but may introduce spurious distinctions ($\eps(\phi)>0$) for
non-faithful ones. Neither recovers what the other destroys in the general case.

\subsection{The Diagnostic Grid}

Each failure mode yields a testable diagnosis:

\begin{enumerate}[leftmargin=1.8em, label=(\arabic*)]
  \item \emph{Projection failure}: $H(X\mid S)$ is large. Distinctions needed
        for the task have been collapsed. Return to a less-projected representation.
  \item \emph{Embedding failure}: $\eps(\phi)>0$ and spurious distinctions
        are driving behavior. Coarsen the embedded space.
  \item \emph{Revision failure}: $\delta_\infty(x)>0$. The residual deficit
        is positive; the system is stuck. Identify what template expansion
        would trigger a productive revision.
  \item \emph{Transport failure}: $\eps(\tau)$ is large relative to the
        oscillation bound in Theorem~\ref{thm:transport-stability}. The analogy
        is invalid; quantify and correct the distortion.
\end{enumerate}

\subsection{The Meta-Theorem}

\begin{theorem}[Meta-theorem on representational adequacy]
\label{thm:meta}
A representational system manages its ontological deficit well iff:
\begin{enumerate}[label=(\roman*)]
  \item \emph{Projection}: $H(X\mid S)$ is small for task-relevant objects.
  \item \emph{Embedding}: $\eps(\phi)$ is controlled; few spurious distinctions.
  \item \emph{Revision}: $\delta_\infty(x)=0$; residual deficit vanishes under
        the reachable revision sequence.
  \item \emph{Transport}: $C\eps(\tau)$ is within acceptable tolerance for
        invariants of interest.
\end{enumerate}
\end{theorem}

Each condition is independently measurable via conditional entropy, false
discovery rate, asymptotic deficit convergence, and transport distortion
respectively.

% ─────────────────────────────────────────────────────────────────────────────
\section{The Deficit as Universal Invariant}
\label{sec:deficit-universal}
% ─────────────────────────────────────────────────────────────────────────────

\subsection{Portability}

\noindent\textbf{Compression.} The deficit is the redundancy of archive $T$
for object $x$. The Revision Monotonicity theorem guarantees that a productive
revision sequence converges to a residual deficit, and the compression jumps
are exactly the productive revisions.

\noindent\textbf{Machine learning.} The deficit is approximation error:
the component of generalization error due to hypothesis class limitation that
more data cannot reduce.

\noindent\textbf{Scientific discovery.} Normal science is the regime of small,
decreasing $\delta_T$ within the current paradigm. Scientific revolution is
a productive revision with large $\Ddelta(\rho)<0$: the paradigm shift reaches
a new $T'$ with much lower deficit for the anomalous observations.

\noindent\textbf{Language evolution.} New terms are productive revisions of
the linguistic archive; each coinage decreases $\delta_T(x)$ for objects
previously requiring circumlocution.

\noindent\textbf{Repair theory.} Repair is transport from a high-deficit
damaged state to a low-deficit repaired state; the validity of the repair
is governed by Transport Stability.

\subsection{Subadditivity}

\begin{theorem}[Deficit Subadditivity]
\label{thm:subadditive}
Let $x\oplus y$ denote an object whose optimal description decomposes
independently into descriptions of $x$ and $y$. Then:
\[
  \delta_T(x\oplus y) \;\le\; \delta_T(x) + \delta_T(y) + O(\log n),
\]
where $n=|x|+|y|$ is the combined description length and the $O(\log n)$
term accounts for the overhead of specifying the decomposition.
\end{theorem}

\begin{proof}
$L_T(x\oplus y)\le L_T(x)+L_T(y)$ since a description of $x\oplus y$
can be formed by concatenating descriptions of $x$ and $y$ reachable in $T$.
By near-additivity of Kolmogorov complexity for independent objects,
$K(x\oplus y)\ge K(x)+K(y)-O(\log n)$.
Hence $\delta_T(x\oplus y)=L_T(x\oplus y)-K(x\oplus y)
\le(L_T(x)+L_T(y))-(K(x)+K(y)-O(\log n))
=\delta_T(x)+\delta_T(y)+O(\log n)$.
\end{proof}

\begin{remark}
This mirrors the subadditivity of Kolmogorov complexity and makes the deficit
behave like a legitimate information quantity. The Deficit Subadditivity Theorem
connects the framework directly to coding theory: the deficit of a composed
system is no worse than the sum of deficits of its parts.
\end{remark}

% ─────────────────────────────────────────────────────────────────────────────
\section{Toward a Categorical Formulation}
\label{sec:category}
% ─────────────────────────────────────────────────────────────────────────────

The categorical formulation must be derived, not imported. We proceed in
strict order: objects, morphisms, operations, functors.

\subsection{The Category $\mathbf{Dist}$}

\begin{definition}[The category $\mathbf{Dist}$]
\label{def:dist-cat}
The category $\mathbf{Dist}$ has:
\begin{itemize}[leftmargin=1.8em]
  \item \emph{Objects}: ontological triples $(X,\sim,T)$.
  \item \emph{Morphisms}: functions $f:(X,\simX,T_X)\to(Y,\simY,T_Y)$
        with associated distortion $\eps(f)\ge 0$ and deficit change
        $\Ddelta(f)\in\mathbb{R}$.
  \item \emph{Identity}: the identity map on $(X,\sim,T)$ with $\eps=0$,
        $\Ddelta=0$.
  \item \emph{Composition}: $(g\circ f)$ has $\eps(g\circ f)\le\eps(f)+\eps(g)$
        (distortions accumulate) and $\Ddelta(g\circ f)=\Ddelta(f)+\Ddelta(g)$
        (deficit changes add).
\end{itemize}
\end{definition}

By including $T$ in the object, the deficit $\delta_T$ is a genuine invariant
of objects, not an additional datum to be tracked separately. This resolves
the under-specification present in earlier formulations.

\subsection{The Quotient Functor and Projection}

\begin{proposition}[Projection Functor]
\label{prop:proj-functor}
Define the quotient functor $Q:\mathbf{Dist}\to\mathbf{Set}$ by
$Q(X,\sim,T)=X/{\sim}$. Every projection $\sigma$ with
$\sim_1\subseteq\sim_2$ induces a quotient map:
\[
  Q(\sigma):X/{\sim_1}\to X/{\sim_2},\quad[x]_{\sim_1}\mapsto[x]_{\sim_2}.
\]
$Q(\sigma)$ is well-defined and surjective. This identifies projection with
quotient formation in $\mathbf{Dist}$.
\end{proposition}

\begin{proof}
Well-definedness: since $\sim_1\subseteq\sim_2$, every $\sim_1$-class is
a subset of a $\sim_2$-class, so the map is consistent. Surjectivity: every
$\sim_2$-class is the image of some $\sim_1$-class.
\end{proof}

\subsection{Morphism Types and the Hierarchy}

The two-level hierarchy maps onto morphism types:

\medskip\noindent\textbf{Projection}: endomorphism with $\Ddelta\ge 0$,
$Q$-image a quotient map.

\medskip\noindent\textbf{Embedding}: morphism into a larger object
$(Y,\sim',T')$ with $Y\supset X$, $\Ddelta\le 0$ (faithful), $\eps\ge 0$
(potential spurious distinctions).

\medskip\noindent\textbf{Revision}: endomorphism changing $T$ while preserving
$X$, $\Ddelta<0$ (productive).

\medskip\noindent\textbf{Transport}: ordinary morphism between distinct objects,
governed by Theorem~\ref{thm:transport-stability}.

\subsection{The Deficit as Functor}

\begin{proposition}[Deficit functor]
\label{prop:deficit-functor}
Define $\delta:\mathbf{Dist}\to\mathbf{Fun}(X,\mathbb{R}_{\ge 0})$ by
$\delta(X,\sim,T)=(x\mapsto\delta_T(x))$. On morphisms, $\delta(f)$ is
the function $\delta_T(x)\mapsto\delta_{T'}(f(x))=\delta_T(x)+\Ddelta(f)$.
This assignment is functorial: $\delta(\mathrm{id})=\mathrm{id}$ and
$\delta(g\circ f)=\delta(g)\circ\delta(f)$.
\end{proposition}

\begin{remark}
This is a well-defined functor because $\delta_T$ depends on $T$, which is
part of the object. The action on morphisms is determined by the signed
deficit change $\Ddelta(f)$, which is part of the morphism data. The
natural transformation interpretation from earlier drafts holds: $\delta$
measures morphism quality, and perfect morphisms are those with $\Ddelta=0$,
$\eps=0$.
\end{remark}

\subsection{Relation to Admissibility Geometry}

The admissibility relation $\mathcal{A}$ specifies the subcategory of
$\mathbf{Dist}$ whose morphisms have $\eps(f)\le\eps_{\max}$ and
$\Ddelta(f)\le\Delta_{\max}$. In the \RSVP{} framework, inadmissible transitions
are those that increase deficit or distortion beyond acceptable thresholds.
The categorical structure of $\mathbf{Dist}$ subsumes admissibility geometry
as the sub-category determined by these bounds, connecting the formal framework
to the broader research programme.

% ─────────────────────────────────────────────────────────────────────────────
\section{Conclusion}
\label{sec:conclusion}
% ─────────────────────────────────────────────────────────────────────────────

Distinguishability geometry proposes the ontological triple $(X,\sim,T)$
as the primitive object for the study of representational change. The deficit
is the invariant. The four operations are the dynamics. The Classification
Conjecture is the central open problem.

\medskip
\noindent\textbf{What is established.}
The Deficit Proxy Theorem makes the coding deficit a computable surrogate for
the geometric distinguishability deficit. The Representational Reduction
Proposition establishes the four operations as the exhaustive qualitative
fiber effects. Each operation has a characteristic theorem.
The Distortion--Deficit Relationship connects the two error measures. The
Meta-Theorem restates representational adequacy as four independently
measurable conditions. The category $\mathbf{Dist}$ is well-defined and the
deficit is a genuine functor.

\medskip
\noindent\textbf{What remains open.}

\begin{enumerate}[leftmargin=1.8em, label=(\arabic*)]
  \item \emph{Full Deficit Correspondence} (Conjecture~\ref{conj:full-corr}):
        does $\Delta_T(x)=0\iff\delta_T(x)=0$?
  \item \emph{Classification Conjecture} (Conjecture~\ref{conj:classification}):
        does every representational transformation factor into a finite
        composition of the four operations?
  \item \emph{Residual deficit}: when does $\delta_\infty(x)=0$? What
        structural property of the revision sequence ensures convergence to
        zero?
  \item \emph{Transport metric}: is there a canonical metric on $\mathbf{Dist}$
        making transport maps Lipschitz?
  \item \emph{Two-level duality}: is there a duality functor
        $\mathbf{Dist}\to\mathbf{Dist}^{\mathrm{op}}$ exchanging projection
        with transport and embedding with revision?
  \item \emph{Unified deficit bound}: does a single expression govern
        $\delta_{T'}(x)\le\delta_T(x)+C(\sigma,\phi,\rho,\tau)$ for arbitrary
        composed transformations?
\end{enumerate}

The strongest reading of the paper is not ``here is a complete theory''
but ``here is a candidate primitive ontology and four conjectured primitive
operations.'' These six open problems define exactly where the theory
consolidates or breaks. They are where the work is.

% ─────────────────────────────────────────────────────────────────────────────
\appendix
% ─────────────────────────────────────────────────────────────────────────────

\section{Formal Definitions: Fiber Bundles and the Quotient Construction}
\label{app:formal}

An ontological triple $(X,\sim,T)$ corresponds to a fiber bundle $E\to B$
where $B=X/{\sim}$ is the base space of equivalence classes, each fiber
$q^{-1}(b)=[x]_\sim$ lies above its class, and the ontology $T$ equips
$E$ with a description system. Projection increases fiber size (bundle
coarsening). Embedding decreases it (bundle refinement). Revision changes
$B$ by changing $T$. Transport maps between distinct bundles.

The deficit $\delta_T(x)$ is an invariant of the triple, not of the bundle
alone: two triples $(X,\sim,T)$ and $(X,\sim,T')$ with the same fiber
structure but different ontologies carry different deficit functions.

\section{The Conservation Law}
\label{app:conservation}

In the event-sourced setting: at $t=0$, $|\Omega_0|=\log|\mathrm{supp}(\mu)|$.
As history $h_t$ is realized, $|H_t|+|\Omega_t|=|\Omega_0|$ because
information transferred to actuality is subtracted from possibility.
The Projection Deficit Bound (Theorem~\ref{thm:proj-deficit}) is consistent:
the $H(X\mid S)$ bits of information destroyed by projection move from
accessible history into the inaccessible remainder $\Omega_t$.

\section{Deficit Examples and the Finite Illustration}
\label{app:deficit}

\noindent\textbf{Compression.} LZ over periodic string: $\Lreach(x)\approx|x|/\log|x|$,
$K(x)\approx\log p+\log(|x|/p)$ for period $p$. Deficit is bits wasted before
the archive learns the period.

\noindent\textbf{\CLIO{} projection error.} For history $h$ with $\sigma(h)=s$
and fiber $F_s$, $\delta_T(h)=H(X\mid S)\ge\log|F_s|$ consistent with
Theorem~\ref{thm:proj-deficit}.

\noindent\textbf{Machine learning.} Linear class over nonlinear data: the
deficit is the approximation error floor. Revision Monotonicity guarantees
that expanding the hypothesis class (a productive revision) decreases this
floor.

\noindent\textbf{Finite illustration revisited.} In Example~\ref{ex:four}
with $X=\{a,b,c,d\}$: the projection $\sigma$ has $|F_p|=2$, so
$\Ddelta(\sigma)\ge\log 2=1$ bit. The revision from $T_0$ to $T_1$ has
$\delta_{T_1}(a)=\delta_{T_1}(b)=0<\delta_{T_0}(a)=\delta_{T_0}(b)$
when working from the projected space: a productive revision. The transport
$\tau:X\to Y$ has $\eps(\tau)=0$ and hence $\Ddelta(\tau)=0$ by
Theorem~\ref{thm:dist-deficit}: a perfect transport.

\end{document}
